The stationary formula can also be obtained as the limit of a finite-horizon calculation. Suppose the same contract is available at each issue date, its borrowing factor is constant, and there are $n$ borrowing dates left before market access ends. Starting from zero terminal borrowing, backward iteration gives
\begin{equation}\label{eq:finite_horizon_recursion}
 b_{H,n}^M=\alpha\sum_{j=1}^{n}\Phi_H^j,
 \qquad d_{H,n}^M=x_H^*(\alpha+b_{H,n-1}^M).
\end{equation}
Each additional date allows future issuance to support current repayment, while the factor $\Phi_H$ discounts that contribution through the same resource and default mechanism. If $\Phi_H<1$, the geometric sum converges to \eqref{eq:stationary}. The convergence supplies an interpretation of the nonexplosive stationary solution; it does not assert that realized debt follows a deterministic path or that default is absent.

Repeated contingency affects every link in this chain. A one-off substitution changes current proceeds against an unchanged future borrowing opportunity. A permanent contractual change also improves the future opportunity used to support current service. The exact amplification in \eqref{eq:stationarygain} accounts for this second effect. It can be large when the borrowing factor is near one, which makes the levels sensitive to the assumed discount rate and growth distribution. The condition $\Phi_H<1$ is consequently substantive: extrapolating the expression through one would not produce a meaningful finite fiscal limit. All reported permanent-relief cases satisfy it.

These comparisons are conditional stationary policy comparisons in an i.i.d. environment. A deterministic warming path, persistent disaster intensity, or a changing surplus rule would generally make the borrowing factor date- or state-dependent. One would then need backward recursion or a larger state-space fixed point instead of a single geometric multiplier. Deferrals require such an enlargement even when growth remains i.i.d., because postponed service is a liability inherited by the next issuer. The scalar multiplier is therefore a useful benchmark for the permanent-relief experiment, rather than a substitute for the inherited-claims calculation.
